\documentclass{article}
\usepackage[utf8]{inputenc}
\usepackage{amsmath}
\usepackage{amsfonts}
\usepackage{amssymb}

\begin{document}
	
	\title{Geometría Diferencial y Ecuaciones de Campo de Einstein}
	\author{Prof. David Duarte}
	\date{\today}
	
	\maketitle
	
	\section*{Enunciado}
	
	Utilizando las herramientas de la geometría diferencial, demuestre cómo las ecuaciones de campo de Einstein describen la curvatura del espacio-tiempo en presencia de materia y energía. Para ello, realice los siguientes pasos:
	
	\subsection*{Introducción Teórica}
	\begin{itemize}
		\item Describa brevemente los conceptos fundamentales de la geometría diferencial necesarios para entender las ecuaciones de campo de Einstein, como variedades diferenciables, métricas y conexiones.
		\item Explique el tensor de curvatura de Riemann y su relación con la métrica del espacio-tiempo.
	\end{itemize}
	
	\subsection*{Ecuaciones de Campo de Einstein}
	\begin{itemize}
		\item Escriba las ecuaciones de campo de Einstein en forma tensorial:
		\begin{equation}
			G_{\mu\nu} + \Lambda g_{\mu\nu} = \frac{8\pi G}{c^4} T_{\mu\nu}
		\end{equation}
		\item Interprete cada término de la ecuación, identificando el tensor de Einstein \( G_{\mu\nu} \), la constante cosmológica \( \Lambda \), la métrica del espacio-tiempo \( g_{\mu\nu} \), y el tensor de energía-momento \( T_{\mu\nu} \).
	\end{itemize}
	
	\subsection*{Aplicación Práctica}
	\begin{itemize}
		\item Seleccione una solución exacta conocida de las ecuaciones de campo de Einstein, como la solución de Schwarzschild o la solución de Kerr.
		\item Calcule explícitamente los componentes del tensor de Einstein para la métrica elegida.
		\item Analice las implicaciones físicas de la solución en términos de la geometría del espacio-tiempo, como la presencia de horizontes de eventos o singularidades.
	\end{itemize}
	
	\subsection*{Discusión}
	\begin{itemize}
		\item Discuta cómo la geometría diferencial proporciona el lenguaje matemático para formular y resolver las ecuaciones de campo de Einstein.
		\item Reflexione sobre la importancia de estas ecuaciones en la comprensión de fenómenos astrofísicos y cosmológicos, como agujeros negros y la expansión del universo.
	\end{itemize}
	
	\section*{Entrega}
	Prepare un informe detallado que incluya todos los cálculos, diagramas y discusiones necesarias para cada uno de los puntos anteriores. El informe debe demostrar comprensión teórica y habilidad para aplicar conceptos matemáticos a problemas físicos.
	
\end{document}
